\documentclass[10pt, letterpaper]{article}
\usepackage[utf8]{inputenc}
\usepackage[T1]{fontenc}
\usepackage{geometry}
\geometry{top=0.75in, bottom=0.75in, left=1in, right=1in}
\usepackage{fancyhdr}
\usepackage{color}
\usepackage{xcolor}
\usepackage{titlesec}
\usepackage{enumitem}
\usepackage{hyperref}
\usepackage{amssymb}
\usepackage[framemethod=TikZ]{mdframed}

% Official Columbia Primary Color (Dark Blue) and Tints
\definecolor{columbiablue_official}{RGB}{0, 33, 71}
\definecolor{columbiablue_light}{RGB}{240, 245, 255}

% Response Box Environment
\newmdenv[
  backgroundcolor=columbiablue_light,
  linecolor=columbiablue_official,
  linewidth=0.5pt,
  leftline=true,
  rightline=false,
  topline=false,
  bottomline=false,
  roundcorner=2pt,
  innerleftmargin=10pt,
  innerrightmargin=10pt,
  innertopmargin=5pt,
  innerbottommargin=5pt,
  skipabove=5pt,
  skipbelow=5pt
]{responsebox}

\pagestyle{fancy}
\fancyhf{}
\renewcommand{\headrulewidth}{0.5pt}
\renewcommand{\headrule}{\hbox to\headwidth{\color{columbiablue_official}\leaders\hrule height \headrulewidth\hfill}}

% Header
\lhead{\small \textbf{\color{columbiablue_official} COLUMBIA UNIVERSITY} \\ Graduate School of Architecture, Planning and Preservation}
\rhead{\small IRB Protocol Submission}

\titleformat{\section}{\large\bfseries\color{columbiablue_official}}{}{0em}{}[\color{columbiablue_official}\titlerule]
\titlespacing*{\section}{0pt}{12pt}{6pt}

\begin{document}

\vspace*{0.2cm}
\begin{center}
    {\Large \textbf{\color{columbiablue_official} Human Subjects Protocol Data Sheet}} \\[0.5em]
    {\large Protocol \# ACYY0820 (Predicting NIMBYism)}
\end{center}

\vspace{0.3cm}

\begin{responsebox}

Columbia University Human Subjects Protocol Data Sheet

\vspace{0.5em}

\vspace{0.5em}

\vspace{0.5em}

\end{responsebox}

\section*{General Information}

\vspace{0.5em}

\vspace{0.5em}

\noindent \textbf{Protocol:}                  ACYY0820(M00Y01)          Protocol Status:           Submitted

\noindent \textbf{Effective Date:}                                      Expiration Date:

\noindent \textbf{Originating Department Code:}                         OAD GSAPP (0704308)

\noindent \textbf{Principal Investigator:}                              Bou Akar, Hiba (hb2541)

From what Columbia campus does this research         Morningside or Lamont Doherty

\noindent \textbf{originate:}

\noindent \textbf{Title:}                     Predicting NIMBYism

Protocol Version \#:                                  Abbreviated Title:         Predicting NIMBYism

\noindent \textbf{Was this protocol previously assigned a number by an IRB:}                       No

\vspace{0.5em}

\vspace{0.5em}

\vspace{0.5em}

Is the purpose of this submission to obtain a "Not Human Subjects Research" determination?

No

\vspace{0.5em}

\vspace{0.5em}

\vspace{0.5em}

\section*{Attributes}

\vspace{0.5em}

\noindent \textbf{Special review type:} Check all that apply or check "None of the Above" box.

\noindent \makebox[0pt][l]{$\square$}\hspace{1.2em}Review for 45 CFR 46.118 Determination (involvement of human subjects is anticipated but is not yet defined)

\noindent \makebox[0pt][l]{$\square$}\hspace{1.2em}Funding review for Administrative IRB approval (such as for Center or Training Grants)

\noindent \makebox[0pt][l]{\color{columbiablue_official}$\boxtimes$}\hspace{1.2em}None of the above

\vspace{0.5em}

\vspace{0.5em}

\noindent \textbf{IRB of record information:} Will a Columbia IRB be the IRB that is responsible for providing review, approval,

and oversight for this study?

Yes

\noindent \textbf{Select the most appropriate response:}

Columbia will be the IRB of record for the study procedures conducted by Columbia researchers (Note:

this response will apply to most submissions).

\vspace{0.5em}

\vspace{0.5em}

Is this research part of a multicenter study?

No

\vspace{0.5em}

\vspace{0.5em}

Please indicate if any of the following University resources are utilized:

\noindent \makebox[0pt][l]{$\square$}\hspace{1.2em} Cancer Center Clinical Protocol Data Management Compliance Core (CPDM)

\noindent \makebox[0pt][l]{$\square$}\hspace{1.2em} CTSA-Irving Institute Clinical Research Resource (CRR)

\noindent \makebox[0pt][l]{$\square$}\hspace{1.2em} CTSA- Irving Institute Columbia Community Partnership for Health (CCPH)

\noindent \makebox[0pt][l]{$\square$}\hspace{1.2em} Columbia University Biobank (CUB)

\noindent \makebox[0pt][l]{\color{columbiablue_official}$\boxtimes$}\hspace{1.2em} None of the above

\vspace{0.5em}

\vspace{0.5em}

\vspace{0.5em}

\section*{Background}

\vspace{0.5em}

\noindent \textbf{Abbreviated Submission:}

The IRB has an abbreviated submission process for multicenter studies supported by industry or NIH

cooperative groups (e.g., ACTG, HVTN, NCI oncology group studies, etc.), and other studies that have a

\vspace{0.5em}

\vspace{0.5em}

complete stand-alone protocol. The process requires completion of all Rascal fields that provide information

regarding local implementation of the study. However, entering study information into all of the relevant Rascal

fields is not required, as the Columbia IRBs will rely on the attached stand-alone (e.g., sponsor's) protocol for

review of the overall objectives.

If you select the Abbreviated Submission checkbox and a section is not covered by the attached stand-alone

protocol, you will need to go back and provide this information in your submission.

\vspace{0.5em}

\vspace{0.5em}

\section*{Study Purpose and Rationale}

Provide pertinent background description with references that are related to the need to conduct this study. If

this is a clinical trial, the background should include both preclinical and clinical data. Be brief and to the point.

\noindent \makebox[0pt][l]{$\square$}\hspace{1.2em} Abbreviated Submission - This information is included in an attached stand-alone protocol. Proceed to the next

\begin{responsebox}

question

This study examines patterns of neighborhood opposition to housing development in Austin,

\end{responsebox}

Texas and evaluates how different stakeholder groups perceive the potential use of predictive or

algorithmic tools in housing policy implementation. Austin has experienced rapid population

growth, rising housing costs, and significant zoning reform, alongside organized opposition from

neighborhood groups and other stakeholders. Existing research on NIMBYism and

homeownership suggests that property ownership, neighborhood characteristics, and local political

context shape opposition to new housing, but there is limited empirical work using detailed

behavioral records (such as protest petitions and public testimony) to characterize these patterns

at the parcel or case level.At the same time, cities and agencies are increasingly interested in

using data and machine learning to anticipate political behavior and manage planning processes,

raising questions about fairness, transparency, and democratic legitimacy. This study will (1) use

existing administrative and public records to model which parcels or owners are most likely to

participate in formal opposition to zoning or development proposals, and (2) use qualitative

interviews to understand how city officials, neighborhood representatives, housing advocates,

developers, and civic technologists view the legitimacy and risks of such predictive tools. The goal

is to inform ethical and democratically accountable use of data-driven methods in housing and land

use governance.

\vspace{0.5em}

\section*{Study Design}

Describe the methodology that will be used in this study, covering such factors as retrospective vs. prospective

data collection, interventional vs. non-interventional, randomized vs. non-randomized, observational,

experimental, ethnography, etc.

\noindent \makebox[0pt][l]{$\square$}\hspace{1.2em} Abbreviated Submission - This information is included in an attached stand-alone protocol. Proceed to the next

question

\begin{responsebox}

This is a mixed-methods, non-interventional, observational study that combines retrospective

\end{responsebox}

analysis of existing records with prospective qualitative data collection.First, the study will conduct

a retrospective record review and secondary analysis of existing administrative and public datasets

from approximately 2018–2025, including zoning and site plan case records, protest or opposition

records, parcel-level appraisal and property data, public meeting minutes and testimony, and

campaign finance records. These data will be linked at the parcel or case level and used to

construct a de-identified analytic dataset for statistical and machine learning analysis of opposition

patterns. No experimental interventions or randomization will be used.Second, the study will

prospectively collect qualitative data through approximately 10 semi-structured interviews with

adult stakeholders (city officials/planners, neighborhood association leaders, housing advocates,

\vspace{0.5em}

developers, and civic technologists). Interviews may be conducted in person, by phone, or via

secure videoconference and may be audio-recorded with consent. Third, the study will include

non-participant observation of public City Council, Planning and Zoning Commission, and

neighborhood association meetings related to housing and zoning, with field notes documenting

how stakeholders discuss and perceive predictive or algorithmic tools. The qualitative and

quantitative components will be analyzed separately and then integrated to provide a more

comprehensive understanding of opposition patterns and stakeholder perspectives.

\vspace{0.5em}

\section*{Statistical Procedures}

Provide sufficient details so that the adequacy of the statistical procedures can be evaluated including power

calculations to justify the number of participants to be enrolled into the study. Definitions of subject terms such

as enrolled and accrued as used for Rascal submissions can be found in the Subjects section.

\noindent \makebox[0pt][l]{$\square$}\hspace{1.2em} Abbreviated Submission - This information is included in an attached stand-alone protocol. Proceed to the next

question

\begin{responsebox}

For the quantitative component, the study will use descriptive statistics and predictive modeling

\end{responsebox}

to identify factors associated with formal opposition to zoning or development proposals at the

parcel or owner level. Candidate predictors will include property characteristics (for example,

assessed value, land use, approximate tenure), spatial variables (for example, proximity to

proposed developments, neighborhood characteristics), and temporal variables (for example, year

of case). Models may include logistic regression and machine learning classifiers (for example,

tree-based methods). Model performance will be evaluated using standard metrics such as

classification accuracy and area under the ROC curve on held-out test sets or cross-validation

folds. The sample size for the quantitative analysis is determined by the number of existing records

available in the defined time period; traditional a priori power calculations are not applicable

because all eligible records in the dataset will be included.For the qualitative component,

approximately 10 interviews are planned, with the final number guided by pragmatic considerations

and thematic saturation (the point at which additional interviews are unlikely to yield substantially

new themes). Interview recordings will be transcribed and analyzed using qualitative thematic

analysis. Transcripts and field notes will be coded to identify recurring themes related to

perceptions of fairness, legitimacy, transparency, and acceptable uses of predictive tools in

housing policy. Quantitative and qualitative findings will then be compared and integrated in the

interpretation phase; formal statistical integration of these components (for example, mixed-

methods meta-inference) will not involve additional hypothesis testing beyond what is described

above.

\vspace{0.5em}

\vspace{0.5em}

\section*{Exempt and Expedited}

\vspace{0.5em}

Is the purpose of this submission to obtain an exemption determination, in accordance with 45CFR46.101(b):

No

\vspace{0.5em}

\vspace{0.5em}

Is the purpose of this submission to seek expedited review , as per the federal categories referenced in

45CFR46.110?

Yes

\begin{responsebox}

Is the risk of harm to which subjects will be exposed as a result of this research no more than minimal?

\vspace{0.5em}

\vspace{0.5em}

\end{responsebox}

Yes

\begin{responsebox}

Select the category or categories of research into which study procedures fall.

\end{responsebox}

\noindent \makebox[0pt][l]{$\square$}\hspace{1.2em} Category 1 - Clinical studies of drugs and medical devices only when condition (a) or (b) is met. (a)

\begin{responsebox}

Research on drugs for which an investigational new drug application (21 CFR Part 312) is not required. (Note:

Research on marketed drugs that significantly increases the risks or decreases the acceptability of the risks

associated with the use of the product is not eligible for expedited review.) (b) Research on medical devices for

which (i) an investigational device exemption application (21 CFR Part 812) is not required; or (ii) the medical

device is cleared/approved for marketing and the medical device is being used in accordance with its

cleared/approved labeling.

\end{responsebox}

\noindent \makebox[0pt][l]{$\square$}\hspace{1.2em} Category 2 - Collection of blood samples by finger stick, heel stick, ear stick, or venipuncture as follows: (a)

\begin{responsebox}

from healthy, nonpregnant adults who weigh at least 110 pounds. For these subjects, the amounts drawn may

not exceed 550 ml in an 8 week period and collection may not occur more frequently than 2 times per week; or

(b) from other adults and children, considering the age, weight, and health of the subjects, the collection

procedure, the amount of blood to be collected, and the frequency with which it will be collected. For these

subjects, the amount drawn may not exceed the lesser of 50 ml or 3 ml per kg in an 8 week period and

collection may not occur more frequently than 2 times per week.

\vspace{0.5em}

\vspace{0.5em}

PLEASE NOTE: If blood is collected through an existing catheter, you do not qualify for expedited review under

this category.

\end{responsebox}

\noindent \makebox[0pt][l]{$\square$}\hspace{1.2em} Category 3 - Prospective collection of biological specimens for research purposes by noninvasive means.

\begin{responsebox}

Examples include: (a) hair and nail clippings in a nondisfiguring manner; (b) deciduous teeth at time of

exfoliation or if routine patient care indicates a need for extraction; (c) permanent teeth if routine patient care

indicates a need for extraction; (d) excreta and external secretions (including sweat); (e) uncannulated saliva

collected either in an unstimulated fashion or stimulated by chewing gumbase or wax or by applying a dilute

citric solution to the tongue; (f) placenta removed at delivery; (g) amniotic fluid obtained at the time of rupture of

the membrane prior to or during labor; (h) supra- and subgingival dental plaque and calculus, provided the

collection procedure is not more invasive than routine prophylactic scaling of the teeth and the process is

accomplished in accordance with accepted prophylactic techniques; (i) mucosal and skin cells collected by

buccal scraping or swab, skin swab, or mouth washings; (j) sputum collected after saline mist nebulization.

\end{responsebox}

\noindent \makebox[0pt][l]{$\square$}\hspace{1.2em} Category 4 - Collection of data through noninvasive procedures (not involving general anesthesia or

\begin{responsebox}

sedation) routinely employed in clinical practice, excluding procedures involving x-rays or microwaves. Where

medical devices are employed, they must be cleared/approved for marketing. (Studies intended to evaluate the

safety and effectiveness of the medical device are not generally eligible for expedited review, including studies

of cleared medical devices for new indications.) Examples include: (a) physical sensors that are applied either

to the surface of the body or at a distance and do not involve input of significant amounts of energy into the

subject or an invasion of the subject=s privacy; (b) weighing or testing sensory acuity; (c) magnetic resonance

imaging; (d) electrocardiography, electroencephalography, thermography, detection of naturally occurring

radioactivity, electroretinography, ultrasound, diagnostic infrared imaging, doppler blood flow, and

echocardiography; (e) moderate exercise, muscular strength testing, body composition assessment, and

flexibility testing where appropriate given the age, weight, and health of the individual.

\end{responsebox}

\noindent \makebox[0pt][l]{\color{columbiablue_official}$\boxtimes$}\hspace{1.2em} Category 5 - Research involving materials (data, documents, records, or specimens) that have been

\begin{responsebox}

collected, or will be collected solely for nonresearch purposes (such as medical treatment or diagnosis).

PLEASE NOTE: If extra tissue is being taken during a routine clinical procedure (i.e. additional tissue that is not

being taken for diagnostic purposes), you do not qualify for expedited review under this category.

\end{responsebox}

\noindent \makebox[0pt][l]{\color{columbiablue_official}$\boxtimes$}\hspace{1.2em} Category 6 - Collection of data from voice, video, digital, or image recordings made for research purposes.

\noindent \makebox[0pt][l]{\color{columbiablue_official}$\boxtimes$}\hspace{1.2em} Category 7 - Research on individual or group characteristics or behavior (including, but not limited to,

\begin{responsebox}

research on perception, cognition, motivation, identity, language, communication, cultural beliefs or practices,

and social behavior) or research employing survey, interview, oral history, focus group, program evaluation,

\vspace{0.5em}

human factors evaluation, or quality assurance methodologies. (NOTE: Some research in this category may be

exempt from the HHS regulations for the protection of human subjects. 45 CFR 46.101(b)(2) and (b)(3). This

listing refers only to research that is not exempt.)

Do all procedures fall into one or more of the categories listed above?

Y

NOTE: This project appears to be eligible for expedited review.

\vspace{0.5em}

\vspace{0.5em}

\vspace{0.5em}

\end{responsebox}

\section*{Funding}

\vspace{0.5em}

Is there any external funding or support that is applied for or awarded, or are you the recipient of a gift, for this

project?

No

\vspace{0.5em}

\vspace{0.5em}

\vspace{0.5em}

\section*{Locations}

{\footnotesize

\begin{verbatim}





Location Type       Facility Name       Domestic or         Geographic          Local IRB Ethics    Local Site

                                        International       Location            Approval            Approval

Columbia/CUMC       Avery Hall









\end{verbatim}

}

\section*{Personnel}

{\footnotesize

\begin{verbatim}





UNI/Phone           Name                Role                Department          Edit/View           Obtaining

                                                                                                    Informed Consent

hb2541              Bou Akar, Hiba      Principal           ARH Instruction     Edit                N

                                        Investigator        (501000X)

212-854-3513

dl3645              Lewis, Daniel       Investigator        ARH Urban           Edit                N

                                                            Planning

713-371-7875                                                (5010302)





\end{verbatim}

}

\section*{Training and COI}

The PI must ensure that each individual that is added as personnel has met the training requirements for this

\noindent \textbf{study (http:}//www.cumc.columbia.edu/dept/irb/education/index.html). For help identifying which research

compliance trainings you may be required to take, visit the Research Compliance Training Finder.

UNI      Name Annual HIPAA HSP             Resear FDA-       S-I     CRC      Good     GCP - GCP        Genetic

\begin{responsebox}

COI              (CITI)   ch with Regulat                    Clinical Third- Refresh Resear

disclos                   Minors ed                          Practic party     er      ch

ure                       (CITI)   Resear                    e (GCP) trackin           Consen

form                               ch                                 g                t

(CITI)

\end{responsebox}

hb2541 Bou       04/10/2          08/24/2

\begin{responsebox}

Akar,   025              023

Hiba

\end{responsebox}

dl3645 Lewis, 11/25/2             11/25/2

\begin{responsebox}

Daniel 025               025

\vspace{0.5em}

\vspace{0.5em}

\vspace{0.5em}

\vspace{0.5em}

\vspace{0.5em}

\end{responsebox}

\section*{Departmental Approvers}

{\footnotesize

\begin{verbatim}





Electronic Signature: Dory Thrasher (501000X) - Other       Date: 11/26/2025

Electronic Signature: Hiba Bou Akar (501000X) - Principal Date: 11/27/2025

Investigator

Electronic Signature: Daniel Lewis (5010302) - Investigator Date: 12/01/2025









\end{verbatim}

}

\section*{Privacy \& Data Security}

\vspace{0.5em}

Indicate the methods by which data/research records will be maintained or stored (select all that apply):

\noindent \makebox[0pt][l]{$\square$}\hspace{1.2em}Hardcopy (i.e., paper)

\noindent \makebox[0pt][l]{\color{columbiablue_official}$\boxtimes$}\hspace{1.2em}Electronic

\begin{responsebox}

Where will the data be stored?

Y

\end{responsebox}

\noindent \makebox[0pt][l]{\color{columbiablue_official}$\boxtimes$}\hspace{1.2em}On a System

\noindent \makebox[0pt][l]{\color{columbiablue_official}$\boxtimes$}\hspace{1.2em}On an Endpoint

\begin{responsebox}

Identify what type of endpoint will be used (select all that apply):

\end{responsebox}

\noindent \makebox[0pt][l]{\color{columbiablue_official}$\boxtimes$}\hspace{1.2em}Desktop Computer

\noindent \makebox[0pt][l]{\color{columbiablue_official}$\boxtimes$}\hspace{1.2em}Laptop Computer

\noindent \makebox[0pt][l]{$\square$}\hspace{1.2em}Mobile Device

\noindent \makebox[0pt][l]{$\square$}\hspace{1.2em}Other

\vspace{0.5em}

\vspace{0.5em}

Does this study involve the receipt or collection of Sensitive Data?

Yes

If any Sensitive Data is lost or stolen as part of your research protocol, you must inform both the IRB and the

appropriate IT Security Office (CUMC IT Security if at CUMC; CUIT if at any other University campus).

\vspace{0.5em}

\vspace{0.5em}

What type of Sensitive Data will be obtained or collected? Select all that apply:

\noindent \makebox[0pt][l]{\color{columbiablue_official}$\boxtimes$}\hspace{1.2em}Personally Identifiable Information (PII), including Social Security Numbers (SSN)

\begin{responsebox}

Will Social Security Numbers (SSNs) be collected for any purpose?

No

\end{responsebox}

\noindent \makebox[0pt][l]{$\square$}\hspace{1.2em}Protected Health Information (PHI), including a Limited Data Set (LDS)

If any PHI is lost or stolen, you must inform both the IRB and the Office of HIPAA Compliance.

\vspace{0.5em}

\vspace{0.5em}

\vspace{0.5em}

Indicate plans for secure storage of electronic sensitive data: check all that apply

\noindent \makebox[0pt][l]{$\square$}\hspace{1.2em}Sensitive data will not be stored in electronic format

\noindent \makebox[0pt][l]{$\square$}\hspace{1.2em}Sensitive data will be stored on a multi-user system

\noindent \makebox[0pt][l]{\color{columbiablue_official}$\boxtimes$}\hspace{1.2em}Sensitive data will be stored on an encrypted endpoint

\vspace{0.5em}

\vspace{0.5em}

\vspace{0.5em}

By Selecting an Endpoint Device and approving this protocol for submission to the IRB, the PI is attesting that

the device and any removable media that may be used have been or will be registered and/or will be

maintained in compliance with the University's Information Security Charter and all related policies. It is

important that this information is updated, during the course of the study, as new devices are added.

\vspace{0.5em}

\vspace{0.5em}

Provide a description of how the confidentiality of study data will be ensured, addressing concerns or

protections that specifically relate to the data storage elements identified above (e.g. hard copy, electronic,

\noindent \textbf{system, and/or endpoint):}

All identifiable data and research records will be stored on password-protected computers and secure, access-controlled

storage managed by Columbia University or other university-approved services. Electronic files will be kept in folders

restricted to the PI and authorized study personnel only. Direct identifiers (such as names, email addresses, and parcel-

level addresses) will be stored in a separate, encrypted linkage file, and the main analytic dataset will use coded study

IDs. Audio recordings and transcripts will be stored in electronic form in the same secure locations and will be labeled by

coded ID rather than names whenever feasible.

\vspace{0.5em}

\vspace{0.5em}

Hardcopy materials (such as signed consent forms and any handwritten notes) will be stored in a locked cabinet in the

PI’s or advisor’s office, with access limited to the research team. Identifiable data will not be sent by unencrypted email or

stored on personal devices that are not password-protected. Only de-identified or aggregate results will be included in

presentations, thesis documents, and publications, and any illustrative quotations from interviews will be edited to remove

names or other details that could reasonably identify individuals.

\vspace{0.5em}

\vspace{0.5em}

If your project is not NIH funded, has a Certificate of Confidentiality (CoC) been requested for this research?

No

\vspace{0.5em}

\vspace{0.5em}

Provide a description of the protections in place to safeguard participants' privacy while information is being

\noindent \textbf{collected:}

Interviews will be scheduled at times and locations that protect participants’ privacy, such as private offices or meeting

rooms, or via secure videoconference or telephone. For remote interviews, a secure, password-protected platform will be

used, and participants will be asked to join from a private space whenever possible. At the start of each interview, the PI

will remind participants that they may decline to answer any question, request that the audio recorder be turned off, or

stop the interview at any time without penalty. Interviews will be conducted one-on-one; no group interviews are planned.

\vspace{0.5em}

\vspace{0.5em}

When audio recording is used, participants will be informed before recording begins and will have the option to participate

without recording if they prefer (in which case only notes will be taken). Field notes from observations of public meetings

will focus on publicly observable behavior and statements that are already part of the public record; covert recording of

private conversations will not occur. Names and specific identifying details will be minimized in notes and transcripts, and

any quotations used in reports will be edited to avoid inadvertently revealing a participant’s identity or employer.

\vspace{0.5em}

\vspace{0.5em}

\vspace{0.5em}

\begin{responsebox}

Procedures

\vspace{0.5em}

\end{responsebox}

Is this project a clinical trial?

No

Is this project associated with, or an extension of, an existing Rascal protocol?

No

Do study procedures involve any of the following?

\begin{responsebox}

Analysis of existing data and/or prospective record review

Yes

Audio, video or photographic recording of research subjects

Yes

Behavioral Intervention?

\vspace{0.5em}

\vspace{0.5em}

No

Biological specimens (collection or use of)

No

Cancer-related research

No

Drugs or Biologics

No

Future use of data and/or specimens

No

Genetic research

No

Home Visits

No

Human Embryonic and/or Human Pluripotent Stem Cells

No

Imaging procedures or radiation

No

Medical Devices

No

Surgical procedures that would not otherwise be conducted or are beyond standard of care

No

\end{responsebox}

Will any of the following qualitative research methods be used?

\begin{responsebox}

Survey/interview/questionnaire

Yes

NOTE: You must attach a PDF version of the survey(s)/interview(s)/questionnaire(s) to this protocol prior

to submission.

Systematic observation of public or group behavior

Yes

Program evaluation

No

\end{responsebox}

Will any of the following tests or evaluations be used?

\begin{responsebox}

Cognitive testing

No

Educational testing

No

Non-invasive physical measurements

No

Taste testing

No

\end{responsebox}

Is this study planning to return any study results to participants?

Yes

\begin{responsebox}

Do plans for returned results include individual participant study results?

No

Do plans for returned results include aggregate study results?

Yes

Please include a plan for return of results in the protocol and, as applicable, information about return of results in the

consent document.

\end{responsebox}

Is there a stand-alone protocol that describes ALL procedures in this study?

\vspace{0.5em}

No

\begin{responsebox}

Please describe ALL study procedures in detail.

NOTE: Be sure to detail all of the procedures above to which a "yes" response was selected. Also detail any

additional procedures that may or may not fall into the categories listed above.

\vspace{0.5em}

\vspace{0.5em}

This study uses two main sets of procedures: (1) secondary analysis of existing records and (2) qualitative interviews

and observations. First, the research team will obtain existing datasets about zoning cases, protest petitions or formal

opposition, parcel-level property characteristics, campaign contributions, and census demographics from public

records, open data portals, and public information requests. These datasets will be linked at the parcel or case level

and used to construct a de-identified analytic dataset for statistical and machine learning analysis of opposition

patterns.

\vspace{0.5em}

\vspace{0.5em}

Second, the study will conduct approximately 10 semi-structured interviews with stakeholders, including city

officials/planners, neighborhood association leaders, housing advocates, developers, and civic technologists.

Interviews may be conducted in person, by phone, or via secure videoconference, with audio-recording (with

participant permission) for transcription and analysis. The study will also involve systematic observation of public City

Council, Planning \& Zoning Commission, and neighborhood association meetings related to housing and zoning over

a six-month period, with field notes documenting how different stakeholders discuss and perceive the use of predictive

or algorithmic tools in housing policy.

\vspace{0.5em}

\vspace{0.5em}

\vspace{0.5em}

Analysis of Existing Data and/or Prospective Record Review

\vspace{0.5em}

\end{responsebox}

Indicate whether the data that will be collected or utilized for the proposed study are in existence as of the

current IRB submission date.

\vspace{0.5em}

\vspace{0.5em}

\begin{responsebox}

Some of the data are in existence and some will be generated in the future.

\vspace{0.5em}

\vspace{0.5em}

Provide the date range of the existing data, documents, or records (e.g., medical charts, school records,

census data)

\vspace{0.5em}

\vspace{0.5em}

Beginning Date: 01/01/2018

End Date:       05/08/2026

\vspace{0.5em}

\vspace{0.5em}

\vspace{0.5em}

As this research involves the collection of some data that will be generated after the IRB protocol is

submitted, informed consent and HIPAA Authorization may be required from subjects

\vspace{0.5em}

\vspace{0.5em}

\vspace{0.5em}

\vspace{0.5em}

\end{responsebox}

\noindent \textbf{Data will be obtained from (select all that apply):}

\vspace{0.5em}

\vspace{0.5em}

\noindent \makebox[0pt][l]{$\square$}\hspace{1.2em} Columbia and/or NYP (e.g., departmental databases/systems, patient charts, Eclipsys, WebCIS,

\begin{responsebox}

administrative/billing records, etc.)

\vspace{0.5em}

\vspace{0.5em}

\end{responsebox}

\noindent \makebox[0pt][l]{\color{columbiablue_official}$\boxtimes$}\hspace{1.2em} Outside Columbia and/or NYP:

\vspace{0.5em}

\vspace{0.5em}

\vspace{0.5em}

\vspace{0.5em}

\begin{responsebox}

Identify the source:

City of Austin open data portal and public records (zoning and site plan cases, council minutes/ATXN

videos), Travis Central Appraisal District parcel/appraisal records, Texas Ethics Commission campaign

finance data, and U.S. Census Bureau demographics

\vspace{0.5em}

\vspace{0.5em}

Based on the types of data to be obtained from outside source(s), additional

documentation/authorization, such as Data Use or Business Associate Agreements, may be

necessary.

\vspace{0.5em}

\vspace{0.5em}

\vspace{0.5em}

\vspace{0.5em}

\end{responsebox}

Will a member of the research team be abstracting data directly from source documents?

\begin{responsebox}

No

\vspace{0.5em}

\vspace{0.5em}

\vspace{0.5em}

\vspace{0.5em}

\end{responsebox}

If any existing data was obtained from a prior research study, was any member of the current research team

involved (e.g., obtained consent, performed study procedures, conducted data analysis) in the project or

procedures that collected and/or used identifiable information?

\vspace{0.5em}

\vspace{0.5em}

\begin{responsebox}

N/A

\vspace{0.5em}

\vspace{0.5em}

\vspace{0.5em}

\vspace{0.5em}

\end{responsebox}

Indicate the manner in which the existing data and/or the records to be reviewed prospectively will be collected

\noindent \textbf{or received:}

(Select all that apply. At least one must be selected.)

\vspace{0.5em}

\vspace{0.5em}

\noindent \makebox[0pt][l]{\color{columbiablue_official}$\boxtimes$}\hspace{1.2em} Contains direct identifiers (e.g., name, MRN, date of birth)

\noindent \makebox[0pt][l]{\color{columbiablue_official}$\boxtimes$}\hspace{1.2em} Coded and the research team has the key and can link the data to direct identifiers

\noindent \makebox[0pt][l]{$\square$}\hspace{1.2em} Coded and the research team does not have access to the key to link data to direct identifiers

\noindent \makebox[0pt][l]{$\square$}\hspace{1.2em} Prior to the receipt of the data by the research team submitting this protocol, the identifiers will be removed and

\begin{responsebox}

no link will remain.

\end{responsebox}

\noindent \makebox[0pt][l]{$\square$}\hspace{1.2em} The information was originally or will be collected without identifiers

\vspace{0.5em}

\vspace{0.5em}

\vspace{0.5em}

If data are collected or received at any point in time with direct identifiers or linked to identifiers, then the data are

considered to be identifiable, and the requirements for Informed Consent (or a waiver, if applicable) and HIPAA

Authorization (or a waiver, if applicable) apply. The necessary information will need to be included in the

respective sections of the submission.

\vspace{0.5em}

\vspace{0.5em}

\vspace{0.5em}

\section*{Recruitment And Consent}

\vspace{0.5em}

\noindent \textbf{Recruitment:}

\vspace{0.5em}

\vspace{0.5em}

Will you obtain information or biospecimens for purposes of screening or determining eligibility?

\vspace{0.5em}

\vspace{0.5em}

\begin{responsebox}

No

\vspace{0.5em}

\vspace{0.5em}

\end{responsebox}

\noindent \textbf{Describe how participants will be recruited:}

Participants will be recruited using purposive and snowball sampling from existing professional networks and publicly

available contact information. The PI will identify potential participants in five stakeholder groups (city officials/planners,

neighborhood association leaders, housing advocates, developers, and civic technologists) and contact them by email or

person-to-person invitation. The recruitment message will briefly describe the study, what participation involves (a 45–60

minute semi-structured interview), and emphasize that participation is voluntary and can be declined or withdrawn at any

time without penalty. Additional participants may be identified through referrals from initial interviewees.

\vspace{0.5em}

\vspace{0.5em}

\noindent \textbf{Select all methods by which participants will be recruited:}

\noindent \makebox[0pt][l]{$\square$}\hspace{1.2em} Study does not involve recruitment procedures

\noindent \makebox[0pt][l]{\color{columbiablue_official}$\boxtimes$}\hspace{1.2em} Person to Person

\noindent \makebox[0pt][l]{$\square$}\hspace{1.2em} Radio

\noindent \makebox[0pt][l]{$\square$}\hspace{1.2em} Newspapers

\noindent \makebox[0pt][l]{$\square$}\hspace{1.2em} Direct Mail

\noindent \makebox[0pt][l]{$\square$}\hspace{1.2em} Website

\noindent \makebox[0pt][l]{\color{columbiablue_official}$\boxtimes$}\hspace{1.2em} Email

\noindent \makebox[0pt][l]{$\square$}\hspace{1.2em} Television

\noindent \makebox[0pt][l]{\color{columbiablue_official}$\boxtimes$}\hspace{1.2em} Telephone

\noindent \makebox[0pt][l]{$\square$}\hspace{1.2em} Flyer/Handout

\noindent \makebox[0pt][l]{$\square$}\hspace{1.2em} Newsletter/Magazine/Journal

\noindent \makebox[0pt][l]{$\square$}\hspace{1.2em} ResearchMatch

\noindent \makebox[0pt][l]{$\square$}\hspace{1.2em} CUMC RecruitMe

\noindent \makebox[0pt][l]{$\square$}\hspace{1.2em} Epic Consent to Contact for Research (CCR) Registry

\vspace{0.5em}

\vspace{0.5em}

\noindent \textbf{Informed Consent Process:}

\vspace{0.5em}

\vspace{0.5em}

\noindent \textbf{Informed Consent Process, Waiver or Exemption:} Select all that apply

\vspace{0.5em}

\vspace{0.5em}

\noindent \makebox[0pt][l]{\color{columbiablue_official}$\boxtimes$}\hspace{1.2em} Informed consent with written documentation will be obtained from the research participant or appropriate

\begin{responsebox}

representative.

\vspace{0.5em}

\vspace{0.5em}

Documentation of informed consent is applicable to:

A portion of the study or subject population

\vspace{0.5em}

\vspace{0.5em}

Identify the portion of the study (e.g., prospective portion, focus groups, substudy 2) or

subject population for which documentation of consent will be obtained::

Semi-structured stakeholder interviews (n=10)

\vspace{0.5em}

\vspace{0.5em}

Documentation of participation will be obtained from::

\end{responsebox}

\noindent \makebox[0pt][l]{\color{columbiablue_official}$\boxtimes$}\hspace{1.2em} Adult participants

\noindent \makebox[0pt][l]{$\square$}\hspace{1.2em} Parent/Guardian providing permission for a child's involvement

\noindent \makebox[0pt][l]{$\square$}\hspace{1.2em} Legally Authorized Representatives (LARs)

\vspace{0.5em}

\vspace{0.5em}

\begin{responsebox}

Describe how participants' written consent will be obtained:

Potential interview participants will receive an email or person-to-person invitation describing the study,

\vspace{0.5em}

\vspace{0.5em}

\vspace{0.5em}

including its purpose, procedures, expected duration, risks, benefits, and voluntary nature. Prior to the

interview, participants will be provided with a written consent document (by email for remote interviews or

in person for in-person interviews) and given time to read it and ask questions. The PI will review the key

elements of the consent form verbally, confirm that the participant understands the study and has had an

opportunity to ask questions, and then obtain the participant’s signature (electronic or handwritten)

before beginning any interview questions. A copy of the signed consent document will be provided to the

participant for their records.

\vspace{0.5em}

\vspace{0.5em}

\end{responsebox}

\noindent \makebox[0pt][l]{$\square$}\hspace{1.2em} Informed consent will be obtained but a waiver of written documentation of consent (i.e., agreement to

\begin{responsebox}

participate in the research without a signature on a consent document) is requested.

\vspace{0.5em}

\vspace{0.5em}

\end{responsebox}

\noindent \makebox[0pt][l]{\color{columbiablue_official}$\boxtimes$}\hspace{1.2em} A waiver of some or all elements of informed consent (45 CFR 46.116) is requested.

\begin{responsebox}

Waiver of consent is applicable to:

A portion of the study or subject population

\vspace{0.5em}

\vspace{0.5em}

Identify the portion of the study (e.g., retrospective chart review portion of the study) or

subject population where a waiver of consent applies:

Existing administrative and public records (zoning, appraisal, council, campaign finance data)

\vspace{0.5em}

\vspace{0.5em}

Select the applicable situation:

\vspace{0.5em}

\vspace{0.5em}

\end{responsebox}

\noindent \makebox[0pt][l]{\color{columbiablue_official}$\boxtimes$}\hspace{1.2em}This study qualifies for a waiver or alteration of consent as the following criteria are met in this study

\begin{responsebox}

(provide justification for EACH of these criteria):

\vspace{0.5em}

\vspace{0.5em}

(1)The research involves no more than minimal risk to the subjects

Provide justification:

The research involves no more than minimal risk to subjects. The existing data consist of public or

administrative records about zoning cases, parcel characteristics, meeting participation, and

campaign contributions involving adult subjects. The primary risk is a potential breach of

confidentiality. This risk will be minimized by storing identifiable data on secure, access-controlled

servers; limiting access to the study team; using coded identifiers in the analytic dataset; and

reporting only aggregate or de-identified results in presentations and publications.

\vspace{0.5em}

\vspace{0.5em}

(2) The waiver or alteration will not adversely affect the rights and welfare of the subjects

Provide justification:

The waiver will not adversely affect the rights or welfare of subjects. The information being used is

drawn from records that are already publicly available or maintained by public agencies (for

example, property tax appraisal records, zoning case files, public meeting testimony, and

campaign finance data). The research does not alter any services, benefits, or legal rights to

which subjects are entitled and does not involve any experimental interventions. Results will be

presented only at aggregate or de-identified levels, so no individual will be singled out in

publications.

\vspace{0.5em}

\vspace{0.5em}

(3) The research could not practicably be carried out without the waiver or alteration

Provide justification:

The research could not practicably be carried out without the waiver. The existing datasets include

potentially thousands of individual property owners, meeting participants, and other stakeholders

across multiple years. Many individuals may no longer reside at the recorded addresses or may

\vspace{0.5em}

be difficult or impossible to contact. Attempting to obtain individual consent from all affected

individuals would be prohibitively burdensome, would substantially delay or prevent the research,

and would likely introduce significant selection bias by including only those who can be contacted

and choose to respond. A waiver is therefore necessary to permit a complete and unbiased

analysis of opposition patterns.

\vspace{0.5em}

\vspace{0.5em}

(4) Whenever appropriate, the subjects will be provided with additional pertinent information

after participation

Provide justification:

Providing additional information after participation is not appropriate or practicable for the records-

based component of this study. The research uses existing records from a large number of

individuals who are not directly contacted or actively enrolled as participants, and many may no

longer be reachable at their recorded addresses. However, the overall findings of the study will be

disseminated through academic and policy publications and presentations, which will be publicly

available and may be accessed by interested individuals and communities.

\vspace{0.5em}

\vspace{0.5em}

(5) If the research involves using identifiable private information or identifiable biospecimens,

the research could not practicably be carried out without using such information or

biospecimens in an identifiable format

\vspace{0.5em}

\vspace{0.5em}

**This waiver criterion does not apply if your research was initially approved before

January 21, 2019, which is the general compliance date for the revised regulations at

45CFR46 subpart A.

\vspace{0.5em}

\vspace{0.5em}

Provide justification:

The research could not practicably be carried out without using identifiable information in an

identifiable format at least during the data linkage and feature-construction stages. Names,

addresses, and other identifiers are necessary to link protest or opposition records to parcel-level

property data, to geocode locations, and to construct derived variables such as owner-occupancy

indicators, approximate length of residence, and neighborhood context measures. Once linkage

and variable construction are complete, these direct identifiers will be removed from the analytic

dataset and replaced with coded study IDs, but some temporary use of identifiable information is

essential to ensure accurate linkage and valid analyses.

\vspace{0.5em}

\vspace{0.5em}

\end{responsebox}

\noindent \makebox[0pt][l]{$\square$}\hspace{1.2em}This study qualifies for waiver or alteration of consent involving public benefit and service programs as

\begin{responsebox}

the following criteria are met for this study (provide justification for EACH of these criteria):

\vspace{0.5em}

\vspace{0.5em}

\end{responsebox}

\noindent \makebox[0pt][l]{$\square$}\hspace{1.2em} Planned Emergency Research with an exception from informed consent as per 21 CFR 50.24.

\vspace{0.5em}

\vspace{0.5em}

\noindent \makebox[0pt][l]{$\square$}\hspace{1.2em} This is exempt research.

\vspace{0.5em}

\vspace{0.5em}

Subject Language

Enrollment of non-English speaking subjects is not expected.

\vspace{0.5em}

\vspace{0.5em}

\begin{responsebox}

During the course of the study, if non-English speaking subjects are encountered, refer to the IRB's

policy on the Enrollment of Non-English Speaking Subjects in Research for further details

(http://www.cumc.columbia.edu/dept/irb/policies/documents/Nonenglishspeakingsubjects.Revised.F

\vspace{0.5em}

\vspace{0.5em}

\vspace{0.5em}

INALDRAFT.111909.website.doc

\vspace{0.5em}

\vspace{0.5em}

\end{responsebox}

\noindent \textbf{Capacity to Provide Consent:}

\vspace{0.5em}

\vspace{0.5em}

\begin{responsebox}

Do you anticipate using surrogate consent or is research being done in an adult population where

capacity to consent may be questionable?

No

\vspace{0.5em}

\vspace{0.5em}

\vspace{0.5em}

\end{responsebox}

\section*{Research Aims \& Abstracts}

\vspace{0.5em}

\noindent \textbf{Research Question(s)/Hypothesis(es):}

\begin{responsebox}

This study is guided by the following primary research question:

How can cities identify and understand patterns of neighborhood opposition to housing

development, and what are the democratic implications of using predictive analytics to

anticipate such opposition?

\end{responsebox}

\noindent \textbf{Sub-questions:}

\begin{responsebox}

What demographic, geographic, and behavioral factors predict whether residents will oppose

specific development projects in Austin?

How do different stakeholder groups (homeowners, renters, housing advocates, city officials,

developers, and civic technologists) perceive the legitimacy and utility of algorithmic tools in

housing policy implementation?

What alternative technological approaches could improve housing development processes

without raising concerns about democratic participation?

\end{responsebox}

The working expectation is that homeownership status, property value, length of residence, and

proximity to proposed developments will be positively associated with opposition, and that

stakeholders will raise concerns about fairness, transparency, and democratic legitimacy when

evaluating predictive tools.

\vspace{0.5em}

\noindent \textbf{Scientific Abstract:}

\begin{responsebox}

This mixed-methods study examines neighborhood opposition to housing development in Austin,

\end{responsebox}

Texas, in the context of recent land use and zoning reforms, including the HOME Initiative and

related changes to parking requirements and minimum lot sizes. The quantitative component uses

parcel- and owner-level data from the Travis Central Appraisal District (2018–2025), City of Austin

site plan and zoning case records, City Council meeting testimony and ATXN recordings,

campaign contribution data from the Texas Ethics Commission, and tract-level census

demographics. These data will be linked and used to construct a set of demographic, geographic,

and behavioral predictors of opposition, such as homeownership status, property value, length of

residence, proximity to proposed developments, and prior participation in protests or testimony.

Supervised machine learning models and geospatial analysis will be used to identify factors

associated with observed opposition and to characterize spatial and temporal patterns from

2018–2025.The qualitative component consists of approximately 10 semi-structured interviews

with key stakeholder groups (city officials/planners, neighborhood association leaders, housing

advocates, developers, and civic technologists), supplemented by observations of public meetings

and document analysis of public comments. Interview transcripts will be analyzed thematically to

understand how different stakeholders interpret the legitimacy, fairness, and democratic

\vspace{0.5em}

\vspace{0.5em}

implications of using predictive models to anticipate opposition, and to identify preferred alternative

uses of technology in the development process. By integrating quantitative modeling of opposition

patterns with qualitative assessments of stakeholder perceptions, the study seeks to (1) clarify the

determinants of NIMBY opposition in Austin, and (2) develop an empirically grounded framework

for evaluating the democratic acceptability and ethical limits of algorithmic tools in housing policy

implementation.

\vspace{0.5em}

\noindent \textbf{Lay Abstract:}

\begin{responsebox}

Austin, Texas is experiencing a severe housing affordability crisis. In response, the city has

\end{responsebox}

recently changed its zoning rules to allow more housing in more places. At the same time, some

neighborhood groups and residents continue to oppose specific development projects through

petitions, public testimony, and other forms of political action. This study aims to understand who is

most likely to oppose new housing, why they do so, and how people feel about the idea of using

computer models to predict where opposition might occur.

The researcher will combine public records about properties, zoning cases, and public meeting

participation with a small number of in-depth interviews with city officials, neighborhood association

leaders, housing advocates, developers, and civic technologists. The data analysis will look for

patterns in where and when opposition happens, and the interviews will explore whether

stakeholders see predictive tools as helpful, fair, or potentially harmful to democratic participation.

The goal is not to design a tool for targeting or discouraging opponents, but to better understand

the risks and possibilities of data-driven approaches in housing policy. The findings are intended to

help Austin and other cities design more transparent, ethical, and inclusive ways to use data and

technology when making decisions about where and how to build new housing.

\vspace{0.5em}

\vspace{0.5em}

\vspace{0.5em}

\section*{Risks, Benefits \& Monitoring}

\vspace{0.5em}

\noindent \textbf{Abbreviated Submission:}

The IRB has an abbreviated submission process for multicenter studies supported by industry or NIH

cooperative groups (e.g., ACTG, HVTN, NCI oncology group studies, etc.), and other studies that have a complete

stand-alone protocol. The process requires completion of all Rascal fields that provide information regarding

local implementation of the study. However, entering study information into all of the relevant Rascal fields is not

required, as the Columbia IRBs will rely on the attached stand-alone (e.g., sponsor's) protocol for review of the

overall objectives. .

If you select the Abbreviated Submission checkbox and a section is not covered by the attached stand-alone

protocol, you will need to go back and provide this information in your submission.

\vspace{0.5em}

\vspace{0.5em}

\section*{Potential Risks}

Provide information regarding all risks to participants that are directly related to participation in this protocol,

including any potential for a breach of confidentiality. Risks associated with any of the items described in the

Procedures section of this submission should be outlined here if they are not captured in a stand-alone protocol.

Risks of procedures that individuals would be exposed to regardless of whether they choose to participate in this

research need not be detailed in this section, unless evaluation of those risks is the focus of this research. When

applicable, the likelihood of certain risks should be explained and data on risks that have been encountered in

past studies should be provided.

\noindent \makebox[0pt][l]{$\square$}\hspace{1.2em} Abbreviated Submission - This information is included in an attached stand-alone protocol. Proceed to the next

question

\vspace{0.5em}

\vspace{0.5em}

\begin{responsebox}

This study involves no more than minimal risk to participants. The primary risk is a potential

\end{responsebox}

breach of confidentiality, particularly because the quantitative component uses administrative and

public records (for example, parcel characteristics, zoning case files, public meeting participation,

and campaign contributions) that could be linked to individuals, and the qualitative component

involves interviews about potentially sensitive political views regarding housing development.For

interview participants, there is also a risk of mild psychological or social discomfort when

discussing politically contentious topics, such as neighborhood opposition, city policy, or

relationships with other stakeholders. There is a small risk that, if an interview quotation were

inadvertently identifiable, participants could experience reputational harm in their professional or

\noindent \textbf{community roles.These risks will be minimized by:} (1) storing identifiable information (names,

contact information, parcel-level addresses) separately from the analytic dataset in encrypted,

access-controlled files; (2) replacing direct identifiers with coded study IDs in all analytic files; (3)

limiting access to identifiable data to the PI and authorized research personnel; (4) storing signed

consent forms and any paper notes in a locked cabinet; and (5) presenting only aggregate or de-

identified results in the thesis and any publications. Interview quotations used in written products

will be edited to remove names, titles, or contextual details that could reasonably identify

individuals or their employers. Participation is voluntary, and interviewees may refuse to answer

any question or stop the interview at any time without penalty.

\vspace{0.5em}

\section*{Potential Benefits}

Provide information regarding any anticipated benefits of participating in this research. There should be a

rational description of why such benefits are expected based on current knowledge. If there is unlikely to be

direct benefit to participants/subjects, describe benefits to society. Please note that elements of participation

such as compensation, access to medical care, receiving study results, etc. are not considered benefits of

research participation.

\noindent \makebox[0pt][l]{$\square$}\hspace{1.2em} Abbreviated Submission - This information is included in an attached stand-alone protocol. Proceed to the next

\begin{responsebox}

question

There is no guaranteed direct benefit to participants. Some interview participants may experience

\end{responsebox}

indirect benefits from having their perspectives heard and reflected in an academic analysis of

housing policy and democratic governance in Austin. They may also appreciate the opportunity to

reflect on the opportunities and risks of data-driven tools in their professional practice or civic

engagement.The primary anticipated benefits are to society and to the scholarly and policy

communities. The study will generate empirical evidence on (1) which factors are associated with

neighborhood opposition to housing development in Austin, and (2) how different stakeholder

groups evaluate the legitimacy and fairness of predictive tools in this context. These findings may

help cities design more transparent, ethical, and democratically accountable approaches to using

data and technology in housing policy implementation. However, these benefits are not guaranteed

and will not directly affect participants’ services, benefits, or legal rights.

\vspace{0.5em}

\section*{Alternatives}

If this research involves an intervention that presents greater than minimal risk to participants, describe available

alternative interventions and provide data to support their efficacy and/or availability. Note, participants always

have the option not to participate in research.

\noindent \makebox[0pt][l]{$\square$}\hspace{1.2em} Abbreviated Submission - This information is included in an attached stand-alone protocol. Proceed to the next

\begin{responsebox}

question

\vspace{0.5em}

\vspace{0.5em}

\vspace{0.5em}

Because this is an observational, minimal-risk social-behavioral study that involves records

\end{responsebox}

analysis and voluntary interviews, the primary alternative is simply not to participate. Individuals

may decline interview invitations without any penalty or loss of benefits to which they are otherwise

entitled. Interview participants may also choose to answer only some questions, decline audio

recording, or withdraw from the study at any time before or during the interview.Outside the

research context, stakeholders remain free to express their views on housing policy and

development through existing public channels (for example, City Council hearings, neighborhood

meetings, or other civic processes) without participating in this study.

\vspace{0.5em}

\section*{Data and Safety Monitoring}

Describe how data and safety will be monitored locally and, if this is a multi-center study, how data and safety

will be monitored across sites as well.

\noindent \makebox[0pt][l]{$\square$}\hspace{1.2em} Abbreviated Submission - This information is included in an attached stand-alone protocol. Proceed to the next

\begin{responsebox}

question

Given that this is a single-site, minimal-risk social-behavioral study with no clinical

\end{responsebox}

interventions, a formal Data and Safety Monitoring Board is not required. Data and safety will be

monitored by the PI under the supervision of the faculty advisor (the Principal Investigator listed on

this protocol).

\noindent \textbf{The PI will be responsible for:} (1) ensuring that informed consent is obtained and documented

appropriately for all interview participants; (2) verifying that data are stored in accordance with the

privacy and security plan described in this submission (for example, encryption, access controls,

separation of identifiers from analytic data); (3) periodically reviewing datasets and documentation

for accuracy, completeness, and proper de-identification; and (4) monitoring for any unanticipated

problems involving risks to participants or others.

If any breach of confidentiality, complaint, or other unanticipated problem occurs, the PI will

promptly report it to the faculty advisor and to the IRB in accordance with Columbia policies and

will take corrective actions (such as revising procedures or enhancing security measures) as

needed. Because the study poses no more than minimal risk and does not involve ongoing

medical or psychological interventions, no stopping rules or interim safety analyses are planned

beyond this ongoing oversight by the PI and advisor.

\vspace{0.5em}

\vspace{0.5em}

\vspace{0.5em}

\section*{Subjects}

\vspace{0.5em}

Unless otherwise noted, the information entered in this section should reflect the number of subjects enrolled or

accrued under the purview of Columbia researchers, whether at Columbia or elsewhere.

\noindent \textbf{Target enrollment:}

10

Number anticipated to be enrolled in the next approval period:

10

Does this study involve screening/assessment procedures to determine subject eligibility?

No

Is this a multi-center study?

No

Does this study have one or more components that apply to a subset of the overall study population (e.g. Phase

1/2, sub-studies)?

No

\vspace{0.5em}

\vspace{0.5em}

\section*{Target Enrollment Demographics}

{\footnotesize

\begin{verbatim}

Population Gender



Females                                  Males                                    Non Specific

0%                                       0%                                       100%



Population Age



0-7                       8-17                   18-65                    >65                      Non Specific

0%                        0%                     100%                     0%                       0%



Population Race



American          Asian            Native Hawaiian Black or African White               More than One     Non-Specific

Indian/Alaskan                     or Other Pacific American                            Race

Native                             Islander

0%                0%               0%               0%              0%                  0%                100%



Population Ethnicity



Hispanic or Latino                       Not Hispanic or Latino                   Non-Specific

0%                                       0%                                       100%



\end{verbatim}

}

\section*{Vulnerable Populations as per 45 CFR 46}

Will children/minors be enrolled

No

Will pregnant women/fetuses/neonates be targeted for enrollment?

No

Will prisoners be targeted for enrollment?

No

\noindent \textbf{Other Vulnerable Populations:}

\noindent \makebox[0pt][l]{$\square$}\hspace{1.2em}Individuals lacking capacity to provide consent

\noindent \makebox[0pt][l]{$\square$}\hspace{1.2em}CU/NYPH Employees/Residents/Fellows/Interns/Students

\noindent \makebox[0pt][l]{$\square$}\hspace{1.2em}Economically disadvantaged

\noindent \makebox[0pt][l]{$\square$}\hspace{1.2em}Educationally disadvantaged

\noindent \makebox[0pt][l]{$\square$}\hspace{1.2em}Non-English speaking

\noindent \makebox[0pt][l]{$\square$}\hspace{1.2em}Other Vulnerable populations

\noindent \makebox[0pt][l]{\color{columbiablue_official}$\boxtimes$}\hspace{1.2em}None of the Populations listed above will be targeted for Enrollment

\noindent \textbf{Subject Population Justification:}

This study will enroll approximately 10 adult stakeholders who are directly involved in or affected by housing and zoning

policy in Austin, including city officials and planners, neighborhood association leaders, housing advocates, developers,

and civic technologists. These individuals are best positioned to describe how they perceive the legitimacy, risks, and

benefits of predictive or algorithmic tools in housing policy implementation. No children, prisoners, or other vulnerable

populations are targeted for enrollment.

Does this study involve compensation or reimbursement to subjects?

No

\vspace{0.5em}

\vspace{0.5em}

\vspace{0.5em}

\section*{Attached Consent Forms}

{\footnotesize

\begin{verbatim}









Number          Copied From           Form Type             Title                   Active/InActive         Initiator

ACYY1682                              Consent               Stakeholder Interview   Active                  Daniel Lewis (dl3645)

                                                            Consent – Predicting

                                                            NIMBYism









\end{verbatim}

}

\section*{Documents}

{\footnotesize

\begin{verbatim}





Archived   Document        Document Type     File Name      Active            Stamped          Date Attached      Created By

           Identifier

No         Predicting_NIMB Information       Predicting_NIMB Y                No               11/25/2025         Daniel Lewis

           Yism_Interview_ Sheet/Verbal      Yism_Interview_                                                      (dl3645)

           Guide           Script            Guide.pdf

No         Predicting_NIMB Recruitment       Predicting_NIMB Y                No               11/25/2025         Daniel Lewis

           Yism_Recruitmen Material          Yism_Recruitmen                                                      (dl3645)

           t_Email                           t_Email.pdf

No         Predicting_NIMB Standalone/Spon   Predicting_NIMB Y                No               11/25/2025         Daniel Lewis

           Yism_Protocol_S sor's Protocol    Yism_Protocol_S                                                      (dl3645)

           ummary                            ummary.pdf









\end{verbatim}

}

\section*{Tasks}

\vspace{0.5em}

\vspace{0.5em}

\vspace{0.5em}

\vspace{0.5em}

\vspace{0.5em}

\vspace{0.5em}
\end{document}